\documentclass[article,paper=b6]{jlreq}
\setlength{\columnseprule}{0.5pt}
\setlength{\columnsep}{5\zw}
\usepackage[left=20mm,right=20mm]{geometry}
\usepackage{amssymb,mathtools,mathcomp}
\usepackage{enumitem,tasks}
\usepackage{physics,physics2}
\settasks{label-width=2\zw}
\usephysicsmodule{ab,ab.braket}
\usepackage{bm}
\usepackage{graphicx,tabularx,here}
\usepackage{tikz,tikz-3dplot}
\usetikzlibrary{math,angles,quotes,calc,intersections}
\usepackage{kyotsutest}
\usepackage{answerbox}
\usepackage[pdftitle={証明},pdfauthor={まさし},hidelinks]{hyperref}
\pagestyle{empty}

\begin{document}
\noindent\textbf{［証明］}

連続する3つの整数は，中央の数を$n$とおくと，
$n − 1$，$n$，$n + 1$ と表せる。

最小の数の平方と最大の数の平方の和から$2$を引いた数は，
\begin{align*}
&\ab(n - 1)^{2} + \ab(n + 1)^{2} - 2\\
{}={} &\ab({n}^{2} - 2n + 1) + \ab({n}^{2} + 2n + 1) - 2\\
{}={} &{n}^{2} - 2n + 1 + {n}^{2} + 2n + 1 - 2\\
{}={} &2{n}^{2}
\end{align*}
である。これは中央の数の平方の$2$倍に等しい。

したがって，3つの連続する整数のうち，
最小の数の平方と最大の数の平方の和から$2$を引いた数は，
中央の数の平方の$2$倍に等しい。

\end{document}